\documentclass[11pt,a4paper]{article}

\usepackage{fontspec}
\setmainfont{times.ttf}[
  Path = {C:/Windows/Fonts/},
  BoldFont = timesbd.ttf,
  ItalicFont = timesi.ttf,
  BoldItalicFont = timesbi.ttf
]
\setmonofont{cour.ttf}[
  Path = {C:/Windows/Fonts/},
  BoldFont = courbd.ttf,
  ItalicFont = couri.ttf,
  BoldItalicFont = courbi.ttf
]
\usepackage[polish]{babel}
\usepackage{geometry}
\usepackage{tabularx}
\usepackage{array}
\usepackage{ragged2e}

\geometry{a4paper, margin=2.3cm}
\setlength{\parindent}{0pt}
\setlength{\parskip}{0.55em}
\linespread{1.05}
\pagestyle{plain}

\newcolumntype{Y}{>{\RaggedRight\arraybackslash}X}
\newcolumntype{L}[1]{>{\RaggedRight\arraybackslash}p{#1}}

\begin{document}

\begin{center}
{\Large\bfseries Agent do odpowiadania na pytania dotyczące trendów i~zależności w~publikacjach naukowych\par}
\vspace{0.7em}
Autorzy: Mateusz Jamroż, Szymon Kiełtyka\par
\vspace{0.25em}
Etap 3. Stan wdrożenia modułów. Spotkanie 19.11.2026.
\end{center}

\vspace{0.4em}

\section*{Stan wdrożenia}

W~tym etapie złożyliśmy moduły, które w~pierwszym dokumencie były jeszcze tylko opisane. Aplikacja działa na komputerze autorów. Użytkownik rozmawia w~oknie, a~przygotowanie zbioru jest osobną pracą: wczytujemy plik, liczymy wektory, przypisujemy kierunki i~składamy bazę wektorową. Pytanie zadane w~oknie nie przelicza wektorów od nowa.

Mamy 621\,515 publikacji w~jednym pliku JSONL o~rozmiarze 1\,979\,261\,045 bajtów. Z~rekordu bierzemy identyfikator, tytuł, rok, DOI, abstrakty, słowa kluczowe i~autorów. Metadane trzymamy w~SQLite, wektory w~Chromie. Na każdą publikację przypada jeden wektor i~przy liczeniu nie wypadł żaden rekord.

Autorów łączymy po ORCID. Gdy go nie ma, bierzemy uuid, a~na końcu identyfikator autora z~rekordu. Jeśli ten sam klucz pojawi się dwa razy przy jednej pracy, zostawiamy pierwsze wystąpienie. Dwie osoby o~tym samym imieniu i~nazwisku, ale z~innym identyfikatorem, zostają osobno. W~profilu widać, która z~nich została dopasowana, razem z~jej pracami.

\section*{Liczenie wektorów}

Cały zbiór wektoryzowaliśmy raz, na Kaggle, na dwóch kartach NVIDIA Tesla T4 po 16\,GB pamięci. Kaggle był potrzebny tylko do tego zadania. Pliki ściągnęliśmy do siebie i~od tego momentu rozmowa działa lokalnie, bez wchodzenia na Kaggle.

Użyliśmy modelu BAAI/bge-m3 z~biblioteki sentence-transformers. Chodziło o~to, żeby polskie pytanie dało się zestawić z~angielskim abstraktem. Model zna oba języki, przyjmuje do 8192 tokenów i~oddaje wektor o~1024 wymiarach. Limit tokenów ustawiliśmy ręcznie, bo bez tego biblioteka potrafi uciąć tekst dużo wcześniej.

Przy pełnym zbiorze odpadły pozostałe modele, które rozważaliśmy wcześniej. multilingual-e5-large dobrze radzi sobie z~krótkim tekstem polskim, ale przy 512 tokenach uciąłby przygotowany opis publikacji. multilingual-e5-base ma ten sam limit. all-mpnet-base-v2 i~all-MiniLM-L6-v2 są nastawione głównie na angielski. W~plikach wektorów i~przy późniejszym pytaniu jest więc BAAI/bge-m3.

\begin{center}
\begin{tabularx}{\textwidth}{|L{0.38\textwidth}|Y|}
\hline
\textbf{Model} & \textbf{Użycie przy pełnym zbiorze} \\
\hline
BAAI/bge-m3 & użyty na Kaggle i~przy pytaniu; 1024 wymiary, do 8192 tokenów \\
\hline
multilingual-e5-large & odrzucony: limit 512 tokenów ucina tekst publikacji \\
\hline
multilingual-e5-base & odrzucony: ten sam limit \\
\hline
all-mpnet-base-v2 & odrzucony: model głównie angielski \\
\hline
all-MiniLM-L6-v2 & odrzucony: model głównie angielski \\
\hline
\end{tabularx}
\end{center}

Do modelu idzie jeden tekst na publikację: tytuł, wszystkie abstrakty, które są w~rekordzie, i~wszystkie słowa kluczowe. Ściągamy zbędne spacje i~ucinamy całość na 3200 znakach. Pracy nie tniemy na kawałki. Wektory normalizujemy, a~w~bazie porównujemy je kosinusem. Nie dokładamy przedrostków w~rodzaju \emph{query} ani \emph{passage}.

Na kartach liczyliśmy w~float16, żeby partia zmieściła się w~pamięci T4, a~na dysk zapisaliśmy float32. Partia miała 16 tekstów na kartę i~szły obie karty naraz. Pytanie użytkownika kodujemy później tym samym modelem, już u~nas, i~też powstaje z~niego jeden wektor.

\begin{center}
\begin{tabularx}{\textwidth}{|L{0.36\textwidth}|Y|}
\hline
\textbf{Parametr} & \textbf{Wartość} \\
\hline
Miejsce obliczenia & Kaggle, dwie karty NVIDIA Tesla T4, po 16\,GB \\
\hline
Model & BAAI/bge-m3, sentence-transformers \\
\hline
Wymiar wektora & 1024 \\
\hline
Porównanie & kosinus, wektory znormalizowane \\
\hline
Długość kontekstu & 8192 tokenów \\
\hline
Tekst jednej publikacji & tytuł, abstrakty i~słowa kluczowe, do 3200 znaków \\
\hline
Wektorów na publikację & 1 \\
\hline
Liczenie / zapis & float16 / float32 \\
\hline
Partia na kartę & 16, obie karty \\
\hline
Rekordy & 621\,515, żaden niepominięty \\
\hline
Czas & 6514 sekund, około godziny i~49 minut \\
\hline
Pliki & 63 pary \texttt{vectors} i~\texttt{ids}; ostatni plik ma 1515 wierszy \\
\hline
\end{tabularx}
\end{center}

Wyszło 63 pary plików, od \texttt{vectors\_00000.npy} do \texttt{vectors\_00062.npy}, i~tyle samo list identyfikatorów. W~pełnym pliku jest 10\,000 wierszy, w~ostatnim 1515. Identyfikator w~pliku tekstowym stoi w~tej samej kolejności co wiersz wektora. Obok leży \texttt{metadata.json} z~nazwą modelu, wymiarem, normalizacją, sposobem porównania, typem liczb, kartami, wielkością partii, limitami, liczbą rekordów i~czasem.

Potem, już na naszym komputerze, wczytaliśmy JSONL do SQLite, na gotowych wektorach przypisaliśmy grupy i~te wektory wstawiliśmy do Chromy. Przy wstawianiu model nie liczył ich drugi raz.

W~kolekcji \texttt{publications} przy wektorze zapisujemy identyfikator publikacji, tytuł, rok, DOI, numer i~nazwę grupy oraz autorów. Identyfikator fragmentu kończy się na \texttt{::0}, bo wektor jest jeden. Ten sam tekst służy później wyszukiwaniu słownemu. Aplikację odpalamy lokalnie: API na porcie 8000 i~okno Streamlit. Przy pierwszym uruchomieniu dociągane są wagi BAAI/bge-m3, jeżeli nie ma ich jeszcze na dysku.

\section*{Klasyfikacja kierunków}

Gotowej nazwy dziedziny w~danych nie ma, więc kierunek rozumiemy jako grupę publikacji o~podobnych wektorach. Na pełnym zbiorze użyliśmy MiniBatchKMeans. Ustawienia, z~którymi ten podział wszedł do bazy, są poniżej.

\begin{center}
\begin{tabularx}{\textwidth}{|L{0.36\textwidth}|Y|}
\hline
\textbf{Parametr} & \textbf{Wartość} \\
\hline
Metoda & MiniBatchKMeans \\
\hline
Liczba grup & 6 \\
\hline
Ziarno & 42 \\
\hline
Wielkość partii & 4096 \\
\hline
Powtórzenia startu & 3 \\
\hline
Dane & zapisane wektory BAAI/bge-m3 \\
\hline
Nazwa grupy & trzy najczęstsze słowa kluczowe \\
\hline
\end{tabularx}
\end{center}

Nazwa jest pomocnicza: bierzemy trzy słowa kluczowe, które w~danej grupie pojawiają się najczęściej, i~składamy z~nich etykietę. Pytania o~kierunki korzystają właśnie z~tego podziału.

\begin{center}
\begin{tabularx}{\textwidth}{|c|Y|r|}
\hline
\textbf{Grupa} & \textbf{Nazwa ze słów kluczowych} & \textbf{Publikacje} \\
\hline
0 & the / Poland / education & 131\,609 \\
\hline
1 & Energy / Biofield / Trivedi & 79\,209 \\
\hline
2 & Poland / the / sustainable development & 149\,497 \\
\hline
3 & mechanical properties / microstructure / and & 88\,590 \\
\hline
4 & Poland / COVID-19 / children & 77\,492 \\
\hline
5 & machine learning / and / the & 95\,118 \\
\hline
\end{tabularx}
\end{center}

Suma liczebności wynosi 621\,515. Czy ten podział dobrze oddaje kierunki, sprawdzimy w~następnym kroku. Sposób sprawdzenia dopiero ustalamy. Jeśli trzeba będzie zmienić liczbę grup albo regułę nazywania, policzymy to z~tych samych wektorów, bez drugiej sesji na Kaggle, i~podmienimy kierunki w~SQLite oraz w~Chromie.

W~kodzie jest też drugi model, TF-IDF z~regresją logistyczną. Miałby odtwarzać etykietę złożoną z~pierwszych słów pierwszego słowa kluczowego tej samej pracy. Na pełnych 621\,515 rekordach go nie uczyliśmy, bo w~zbiorze nie ma oficjalnych nazw dziedzin. Odpowiedzi o~kierunkach biorą się z~grup na wektorach, nie z~tego klasyfikatora.

\begin{center}
\begin{tabularx}{\textwidth}{|L{0.36\textwidth}|Y|}
\hline
\textbf{Parametr} & \textbf{Wartość} \\
\hline
Cechy & TF-IDF, do 5000 cech, pojedyncze słowa i~pary \\
\hline
Klasyfikator & regresja logistyczna \\
\hline
Liczba przejść & 1000 \\
\hline
Wagi klas & wyrównane \\
\hline
Część testowa w~kodzie & 30 procent \\
\hline
Ziarno podziału & 42 \\
\hline
Pełny zbiór & nieuczony \\
\hline
\end{tabularx}
\end{center}

\section*{Wyszukiwanie, zestawienia i~odpowiedź}

Pytanie o~treść zamieniamy na wektor modelem BAAI/bge-m3 i~szukamy najbliższych publikacji. Obok działa wyszukiwanie słowne BM25, które lepiej trzyma dokładny tytuł, nazwisko i~rzadki termin. Oba wyniki łączymy metodą RRF, bez porównywania punktacji z~dwóch różnych skal. Domyślnie zwracamy 10 publikacji. W~oknie tę liczbę można ustawić od 1 do 15 i~dodać filtr roku. Pod odpowiedzią podajemy tytuł, rok i~DOI, jeśli jest w~rekordzie.

Liczebności liczymy w~SQLite, bo sam wektor nie powie, ile prac z~danego kierunku wypadło w~kolejnych latach. Do trendu bierzemy lata z~zakresu 1900--2026. Rok spoza tego zakresu zostaje w~bazie, ale nie wchodzi do szeregu. Wzorzec rozwoju porównuje pierwszy i~ostatni punkt: zmiana powyżej 10\% to wzrost, spadek o~więcej niż 10\% to spadek, a~pomiędzy tym poziom zostaje podobny. Gdy w~szeregu jest tylko jeden rok, podajemy liczebność za ten rok i~nie nazywamy tego wieloletnim trendem. Osobno aplikacja pokazuje najliczniejsze kierunki, profil autora, współautorów i~autorów danego kierunku. Jak osoby nie ma w~zbiorze, mówimy o~tym wprost.

Zdanie układa lokalny model llama3.2:3b, uruchomiony na tym samym komputerze. Dostaje już policzone wiersze i~znalezione publikacje, więc nie wymyśla tytułów, lat ani osób. Odpowiedź jest w~języku pytania. Gdy model nie odpowie, zdanie składamy wprost z~tych samych wyników.

\begin{center}
\begin{tabularx}{\textwidth}{|L{0.36\textwidth}|Y|}
\hline
\textbf{Parametr} & \textbf{Wartość} \\
\hline
Model & llama3.2:3b \\
\hline
Dostęp & lokalne API na komputerze autorów \\
\hline
Temperatura & 0,2 \\
\hline
Kontekst & 2048 tokenów \\
\hline
Długość odpowiedzi & do 512 tokenów \\
\hline
Czas oczekiwania & 90 sekund \\
\hline
Gdy model nie odpowie & zdanie z~wyników wyszukiwania i~zestawień \\
\hline
\end{tabularx}
\end{center}

Użytkownik może mieć kilka czatów i~wracać do każdego z~nich. Przy nowym czacie zaznacza, czy ma on dostać historię innego czatu. Dopytanie w~stylu „a kto to napisał” opiera się na pracy albo osobie z~tej historii. Jak nie da się jej rozstrzygnąć, prosimy o~doprecyzowanie. Pusta wypowiedź, powitanie i~pytanie spoza publikacji nie ruszają bazy. Przy liczebności kierunków i~szeregu lat w~oknie jest też wykres.

Kod sprawdzamy lokalnie, bez kontenerów. Testy obejmują klucz autora, złożenie tekstu z~polskiego i~angielskiego abstraktu, rozpoznanie rodzaju pytania, opis wzrostu i~spadku oraz odcięcie lat 1900--2026. Mamy też plik z~pytaniami na małym zbiorze próbnym. Ocenę wyszukiwania na pełnym indeksie i~ocenę podziału na grupy zostawiamy na później, bo sposób tej drugiej dopiero ustalamy.

\section*{Co jeszcze dodajemy}

Wektory BAAI/bge-m3 z~Kaggle zostają. Nie liczymy ich drugi raz. Zostały rzeczy, które były w~planie i~których jeszcze nie domknęliśmy.

Najpierw chcemy ustalić, jak w~ogóle poznamy, że podział na grupy jest sensowny. Dopiero wtedy zdecydujemy, czy zostaje sześć grup i~obecne nazwy, czy zmieniamy liczbę grup albo sposób układania nazwy ze słów kluczowych.

Do pełnego zbioru przygotujemy pytania po polsku i~po angielsku, z~wskazanymi publikacjami, które powinny się pojawić. Porównamy samo wyszukiwanie słowne, samo wyszukiwanie wektorem BAAI/bge-m3 i~połączenie, które jest już w~aplikacji. Osobno wejdą pytania polskie do angielskiego abstraktu, bo pod ten przypadek braliśmy ten model.

Liczby z~odpowiedzi zestawimy z~zapytaniem puszczonym wprost na danych: liczba prac w~kierunku, ranking kierunków i~lata w~szeregu. Do układania zdań podłączony jest teraz llama3.2:3b. Na tych samych pytaniach porównamy go z~Qwen3 4B i~Qwen3 8B i~zostawimy najmniejszy model, który trzyma się przekazanych rekordów i~da się go normalnie używać u~nas.

Na pełnym indeksie zmierzymy też czas startu, pamięć i~czas jednego pytania o~treść, o~trend i~o~autora. Pomiar zrobimy na komputerze autorów, bez kart z~Kaggle. Chcemy wiedzieć, czy rozmowę da się prowadzić bez długiej przerwy między zdaniami.

\begin{center}
\begin{tabularx}{\textwidth}{|L{0.36\textwidth}|L{0.20\textwidth}|Y|}
\hline
\textbf{Zadanie} & \textbf{Stan} & \textbf{Uwaga} \\
\hline
Wczytanie pełnego pliku & ukończone & 621\,515 rekordów w~SQLite \\
\hline
Wektory BAAI/bge-m3 na Kaggle & ukończone & dwie karty T4, 6514 sekund \\
\hline
Baza Chroma & ukończona & gotowe wektory, bez ponownego liczenia \\
\hline
Podział na sześć grup & wdrożony & sposób oceny dopiero ustalamy \\
\hline
Wyszukiwanie i~okno rozmowy & wdrożone & porównanie trzech wariantów przed nami \\
\hline
Porównanie modeli językowych & przed nami & teraz llama3.2:3b; obok Qwen3 4B i~8B \\
\hline
Czas i~pamięć pełnego zbioru & przed nami & pomiar na komputerze autorów \\
\hline
\end{tabularx}
\end{center}

\section*{Przypadki użycia}

Przypadki z~pierwszego dokumentu da się już przejść na pełnym zbiorze. Poniżej jest, jak to wygląda teraz.

\begin{center}
\begin{tabularx}{\textwidth}{|c|L{0.28\textwidth}|Y|}
\hline
\textbf{ID} & \textbf{Przypadek} & \textbf{Stan} \\
\hline
PU-01 & Wczytanie zbioru &
W~SQLite i~w~Chromie jest 621\,515 publikacji. Brak pliku nie kasuje tego, co już zostało wczytane. \\
\hline
PU-02 & Znalezienie prac po treści &
Pytanie idzie wektorem BAAI/bge-m3 i~wyszukiwaniem słownym. Porównanie wariantów na pełnym indeksie jest przed nami. \\
\hline
PU-03 & Dominujące kierunki &
Odpowiedź bierze sześć zapisanych grup. Jak ocenimy ten podział, ustalimy później. \\
\hline
PU-04 & Wzorzec rozwoju kierunku &
Lata 1900--2026. Przy jednym roku podajemy liczebność za ten rok. \\
\hline
PU-05 & Autorzy wskazanej pracy &
Pokazujemy osoby zapisane przy tym rekordzie. \\
\hline
PU-06 & Dorobek osoby &
Widać dopasowane osoby, ich publikacje, lata i~kierunki. Brak osoby mówimy wprost. \\
\hline
PU-07 & Autorzy kierunku &
Bierzemy osoby, których prace wpadły do danej grupy, razem z~tymi pracami. \\
\hline
PU-08 & Dopytanie &
Kolejne zdanie w~czacie może dotyczyć pracy albo osoby z~historii. Jak się nie da jej ustalić, prosimy o~doprecyzowanie. \\
\hline
PU-09 & Wiele czatów &
Nowy czat może dostać historię innego czatu albo zacząć bez niej. Do starych czatów da się wrócić. \\
\hline
PU-10 & Wypowiedź pusta albo spoza zbioru &
Baza zostaje bez zmian. Przy powitaniu proponujemy pytania o~publikacje. \\
\hline
\end{tabularx}
\end{center}

\vspace{1.2em}
Podpis promotora (akceptacja etapu 3): \dotfill

\vspace{0.8em}
Data: \dotfill

\end{document}
